% Authored. Numbers come from generated/macros.tex.

\section{Change-Point Sensitivity in Full}
\label{sec:changepoints}

Section~\ref{sec:robustness} summarises the sensitivity of the piecewise-constant
mean shifts. This appendix records the specification and the full grid, so the
summary can be checked rather than taken.

\subsection{Specification}

For a balance ratio $y_1,\dots,y_n$ inside one statistical regime, the model
partitions the index range into $k+1$ contiguous segments and fits a constant to
each. The partition minimises total within-segment squared error,
\[
\min_{0 = b_0 < b_1 < \cdots < b_k < b_{k+1} = n}
\ \sum_{j=0}^{k} \sum_{t = b_j + 1}^{b_{j+1}} \left(y_t - \bar{y}_{[b_j, b_{j+1}]}\right)^2,
\]
solved exactly by dynamic programming rather than by a greedy or binary-segmentation
heuristic. Selection over $k$ uses
\[
\mathrm{BIC}(k) = n \log\!\left(\frac{\mathrm{RSS}(k)}{n}\right) + (2k+1)\log n,
\]
where the parameter count $2k+1$ charges for $k+1$ segment means and $k$ break
locations. Charging for the locations is the conservative choice: it penalises an
additional break twice and makes the model less willing to claim one. Zero breaks is
always admissible.

Two tuning parameters are not estimated: the minimum segment length and the maximum
number of breaks. The baseline sets them to five years and two breaks per regime.

\subsection{The grid}

Detection is re-run over the product of

\[
\text{minimum segment} \in \{4, 5, 6, 7\}
\quad \text{and} \quad
\text{maximum breaks} \in \{1, 2, 3\},
\]

giving \BreakSpecifications\ specifications for each of the \BreakSeries\
regime-subsector series. Three quantities are persisted for each series: the modal
number of breaks and the share of the grid selecting it, the modal set of dates and
the share of the grid returning exactly that set, and the number of distinct date
sets the grid produced.

The last of these is the most informative single number. A series whose grid returns
one date set is telling us something about the series; a series whose grid returns
six is telling us about the tuning parameters.

\subsection{Reading the result}

\BreakModalAgree\ of \BreakSeries\ series have a preferred specification that agrees
with the modal grid outcome. Where the two disagree, the modal column is the more
reliable summary and the preferred date should not be quoted alone. The least stable
series is \WeakestBreakSector\ in the modern regime, at
\WeakestBreakShare\% modal agreement across \WeakestBreakSets\ distinct date sets.

Two further limits apply regardless of the grid. A minimum segment length of five
years means a shift within the last four years of the sample is undetectable by
construction. And selecting a mean shift does not imply the series is
piecewise-constant; it is the best fit inside a restricted model class, chosen
because the sample is too short to support a richer one.

The full per-specification grid, the BIC ladder over admissible break counts and the
stability summary are persisted as
\path{outputs/tables/structural_break_sensitivity.csv},
\path{outputs/tables/structural_break_bic_ladder.csv} and
\path{outputs/tables/structural_break_stability.csv}.
